\documentclass[12pt,a4paper,finnish]{vendor/tutthesis}

% LaTeX-tiedosto opinnäytepohjaksi
% Tekijä: Sami Paavilainen
% Muokkaaja: Heikki Huttunen (heikki.huttunen@tut.fi) 31.7.2012.

% Vaatii lisäksi luokkatiedoston tutthesis.cls
% sekä tiedoston tty-logo.pdf (pdflatex) tai tty-logo-eps (latex)

\usepackage[finnish]{babel} % tavutus suomeksi, vaihda jos eri kieli
\usepackage[small]{caption}% kuvatekstien kirjasinkoko ja asettelu
%omia käskyjä voi lisätä tähän väliin, esim. \newcommand{\angs}{\textsl{\AA}}
\pagenumbering{Roman}
\graphicspath{ {./vendor/images/} {./images} }
% alkuvalmistelut loppuu

\begin{document}


\thispagestyle{empty}

\vspace*{-.5cm}\noindent

% If thesis is in English, use the file "tut-logo"
% instead of "tty-logo" in the following:

\includegraphics[width=8cm]{tty-logo}

\vspace{6.8cm}

\noindent{\bf\large \textsf{TEKIJÄN NIMI}}\\
{\bf\large \textsf{OPINNÄYTETYÖN OTSIKKO}}\\
\textsf{Diplomityö}

\vspace{8.7cm} % jos kaksi otsikkoriviä vaihda -> 6.7cm

\begin{flushright}

\begin{minipage}[c]{6.8cm}
\begin{spacing}{1.0}
\textsf{Tarkastaja: Tarkastaja 1}\\
\textsf{Tarkastaja ja aihe hyväksytty}\\
\textsf{xxxxxxx tiedekuntaneuvoston}\\
\textsf{kokouksessa xx.xx.xxxx}\\
\end{spacing}
\end{minipage}
\end{flushright}

\newpage

\setcounter{page}{1} % tämä tarvitaan, jottei ensimmäinen sivu kansilehden jälkeen olisi numero 2.

\chapter*{TIIVISTELMÄ}
\begin{spacing}{1.0}
\textsf{TAMPEREEN TEKNILLINEN YLIOPISTO}\\
\textsf{xxxxxxxxxxxxx koulutusohjelma}\\
{\bf \textsf{TEKIJÄN NIMI: Otsikko}}\\
\textsf{Diplomityö, xx sivua, x liitesivua}\\
\textsf{Xxxxxkuu 201x}\\
\textsf{Pääaine: }\\
\textsf{Tarkastajat: }\\
\textsf{Avainsanat: }\\
\end{spacing}

\noindent
Ensimmäinen kappale

\noindent
Toinen kappale
\newpage
\chapter*{ABSTRACT}
\begin{spacing}{1.0}
\textsf{TAMPERE UNIVERSITY OF TECHNOLOGY}\\
\textsf{Master's Degree Programme in xxxxxxx Technology}\\
{\bf \textsf{AUTHOR : Title}}\\
\textsf{Master of Science Thesis, xx pages, x Appendix pages}\\
\textsf{xxxxxx 201x}\\
\textsf{Major: }\\
\textsf{Examiner: }\\
\textsf{Keywords: }\\
\end{spacing}

\noindent
First paragraph

\noindent
Second paragraph

\newpage

\chapter*{ALKUSANAT}
\noindent Tämä (\textit{d-tyo.tex}) on \LaTeX-pohja Tampereen teknillisen
yliopiston opinnäytetöitä varten. Samaan pakettiin kuuluu myös
tiedosto \mbox{\textit{tutthesis.cls}}, joka sisältää taittoteknisiä
lisäyksiä \LaTeX:n alkuperäiseen \textit{report.cls}-luokkatiedostoon.

Lisäksi otsikkosivua varten tarvitaan tiedosto \textit{tty-logo.xxx}, jonka
tulee sisältää TTY:n logo. Tiedoston tulee
olla joko \textit{.eps}- tai \textit{.pdf}-muodossa riippuen \LaTeX-versiosta.

\newpage
\renewcommand\contentsname{Sisällys} % poista jos kieli muu kuin suomi
\tableofcontents
%\newpage
%\listoffigures
%\listoftables
\newpage
\chapter*{TERMIT JA NIIDEN MÄÄRITELMÄT}


\begin{termiluettelo}

\item [$\hbar$] Redusoitu Planckin vakio
\item [SNR] Signaali-kohinasuhde (engl.: Signal to Noise Ratio)

\end{termiluettelo}


\newpage
\renewcommand{\chaptermark}[1]{\markboth{\thechapter. \ #1}{}}
\renewcommand{\sectionmark}[1]{\markright{}{}}
\lhead{\fancyplain{}{\leftmark}}

% Tästä alkaa varsinainen teksti
\pagenumbering{arabic}

\chapter{Johdanto}

Tampereen teknillisen yliopiston perustutkintojen opinnäytetöitä ovat
kandidaatintyö ja diplomityö. Lisäksi tehdään kirjallisia
harjoitustöitä, joissa tätä ohjetta voidaan käyttää soveltuvin osin.

Tämä pohja on tarkoitettu auttamaan työn kirjoittamista \LaTeX:n
kanssa. Oppaita itse \LaTeX:n käyttöön löytyy kattavasti verkosta,
esimerkiksi Wikipediasta. Myös seuraava opas on hyvä lähtökohta
aloittelijoille:
\begin{quote}
http://www.ctan.org/tex-archive/info/lshort/english/lshort.pdf
\end{quote}

\chapter{Opinnäytetyön rakenne}

Tässä esimerkkejä luvuista ja aliluvuista.

\section{Lähdeluettelo}

Lähdeluettelo voi \LaTeX:ssa tehdä joko tietokantamuotoisesti
\textit{bibtex}-käskyn avulla tai vaihtoehtoisesti tässä esitetyllä tyylillä.
Lisätietoa löytyy esim. Wikipediasta ''\textit{bibtex}''-hakusanalla.

Tältä näyttää viittaus perinteisellä numeromerkinnällä \cite{Hirs,Mittelbach}.

% Huom.: oikeaoppinen tapa viitata olisi seuraava:
%
% ...numeromerkinnällä~\cite{Hirs}.
%
% Merkintä "~" toimii kuten välilyönti, mutta estää rivinvaihdon
% ja seuraavan rivin alkamisen rumasti viittauksella.

\subsection{Aliesimerkki}

Turhaa tekstiä esimerkin vuoksi.

\section{Matemaattisista merkinnöistä}

\LaTeX:lla yhtälöiden kirjoittaminen on helppoa. Seuraavassa ensin esimerkki
\begin{equation}
  \label{eq:fourier}
  G^+(t,t')= \int G^+(E) exp[-iE(t-t')/\hbar] dE,
\end{equation}
jonka jälkeen kappale jatkuu näin.
Tekstissä kaavat $G^+(t,t')= \int G^+(E) exp[-iE(t-t')/\hbar] dE$ näyttävät
tällaisilta. Kaavoihin (ks. \ref{eq:fourier}) on myös helppo viitata.

% Esimerkkikuva.
% Huomaa, että useimmat ohjelmat latovat PDF-kuvan
% A4-sivulle, jolloin kuvassa on mukana paljon tyhjää.
% Tämän voi poistaa PDF-editoreilla. Esim.
% PDF-X-Change-ohjelmistossa on toiminto "remove white spaces".
%
% Huomaa myös, että Latex sijoittaa kuvan aikaisintaan
% tähän kohtaan tekstissä. Yleensä kuvat kannattaa
% ladata n. sivun verran aikaisemmin kun niihin viitataan.
% Muuten seuraava sopiva paikka saattaa tulla vasta
% parin sivun päähän viittauksesta.

\begin{figure}[t]
\begin{center}
\end{center}
\caption{Matlabilla tehty PDF-muotoinen esimerkkikuva.}
\label{fig:esimKuva}
\end{figure}

\section{Kuvat}

Latexiin voi tuoda kuvia useassa eri formaatissa.
Näppärintä on käyttää pdflatex-kääntäjää, jolloin
sallitut formaatit ovat \verb+*.png+, \verb+*.jpg+
sekä \verb+*.pdf+. Kuvassa~\ref{fig:esimKuva} on
PDF-formaatissa tuotu Matlabilla tehty
esimerkkikuva.

% lähdeluettelo
\newpage
\addcontentsline{toc}{chapter}{Lähteet} % Lisää tämäkin sisällysluetteloon
\renewcommand{\bibname}{Lähteet}
\begin{thebibliography}{99}
\bibitem{Hirs} Hirsjärvi, S., Remes, P., ja Sajavaara, P. 2005. Tutki ja
  kirjoita, 11. painos. Helsinki, Tammi. 436 s.
\bibitem{Mittelbach} Mittelbach, F., Goossens, M., Braams, J.,
Carlisle, D., Rowley, C. 2004. The Latex Companion, 2. painos. Boston, Addison-Wesley. 1120 s.

\end{thebibliography}

% seuraavia voi käyttää bibtex:n kanssa (edellyttää tiedoston kirjat.bib)
%\bibliographystyle{plain}
%\bibliography{kirjat}

\appendix
\renewcommand{\chaptername}{Liite}

\chapter{Liitteitä}
\end{document}

